\documentclass[12pt,a4paper,reqno]{amsart}

\usepackage[lmargin=1.4 in, rmargin= 1.4in, tmargin=1.7in]{geometry}

\title{The category of necklaces is a test category}

\author{Arne Mertens}
\address[Arne Mertens]{Universiteit Antwerpen, Departement Wiskunde, Middelheimcampus,
Middelheimlaan 1, 2020 Antwerp, Belgium}
\email{arne.mertens@uantwerpen.be}

\usepackage{graphicx}
\usepackage[english]{babel}
\usepackage{amsmath}
\usepackage{amssymb}
\usepackage{amsthm}
\usepackage[all,cmtip]{xy}
\usepackage{cancel}
\usepackage[alpine]{ifsym}
\usepackage{enumerate}
\usepackage{stmaryrd}
\usepackage{tikz}
\usepackage{tikz-cd}
\usepackage{quiver}
\usetikzlibrary{matrix}
\usepackage[toc,page]{appendix}
\usepackage{hyperref}
\usepackage{mathtools}
\usepackage{float}
\usepackage{hhline}
\usepackage[style=alphabetic]{biblatex}
\addbibresource{bibfile.bib}
\setcounter{biburlnumpenalty}{9999}
\setcounter{biburllcpenalty}{9999}
\setcounter{biburlucpenalty}{9999}

%\setcounter{tocdepth}{2}


%Basic operators
\DeclareMathOperator{\sgn}{sgn}
\DeclareMathOperator{\cvee}{\vee\cdots\vee}

%Basic categorical operators
\DeclareMathOperator{\Ob}{Ob}
\DeclareMathOperator{\Mor}{Mor}
\DeclareMathOperator{\Hom}{Hom}
\DeclareMathOperator{\id}{id}

\DeclareMathOperator{\Fun}{Fun}
\DeclareMathOperator{\Nat}{Nat}

\DeclareMathOperator{\Lan}{Lan}
\DeclareMathOperator{\Ran}{Ran}
\DeclareMathOperator*{\colim}{colim}
\DeclareMathOperator{\im}{im}


%Special categories
\DeclareMathOperator{\Set}{Set}
\DeclareMathOperator{\Mon}{Mon}
\DeclareMathOperator{\Comon}{Comon}
\DeclareMathOperator{\Ab}{Ab}
\DeclareMathOperator{\Mod}{Mod}
\DeclareMathOperator{\Vect}{Vect}
\DeclareMathOperator{\Top}{Top}
\DeclareMathOperator{\Ch}{Ch}
\DeclareMathOperator{\Alg}{Alg}
\DeclareMathOperator{\Coalg}{Coalg}
\DeclareMathOperator{\Poset}{Poset}
\DeclareMathOperator{\Lin}{Lin}
\DeclareMathOperator{\SSet}{SSet}
\DeclareMathOperator{\CSet}{CSet}
\DeclareMathOperator{\Lax}{Lax}
\DeclareMathOperator{\Colax}{Colax}
\DeclareMathOperator{\StrMon}{StrMon}
\DeclareMathOperator{\Frob}{Frob}

\DeclareMathOperator{\Quiv}{Quiv}
\DeclareMathOperator{\Cat}{Cat}
\DeclareMathOperator{\MonCat}{\underline{MonCat}}
\DeclareMathOperator{\QCat}{QCat}


%Simplicial set operators
\DeclareMathOperator{\sk}{sk}
\DeclareMathOperator{\cosk}{cosk}

\DeclareMathOperator{\InnHorn}{InnHorn}
\DeclareMathOperator{\LHorn}{LHorn}
\DeclareMathOperator{\RHorn}{RHorn}
\DeclareMathOperator{\Horn}{Horn}
\DeclareMathOperator{\Cell}{Cell}


%Commands
\newcommand{\fint}{\mathbf{\Delta}_{f}} %category of finite intervals
\newcommand{\simp}{\mathbf{\Delta}} %simplicial walking category
\newcommand{\asimp}{\mathbf{\Delta}_{+}} %augmented simplicial walking category
\newcommand{\nsimp}{\mathbf{\Delta}_{-}} %narrow simplicial walking category
\newcommand{\nec}{\mathcal{N}ec} %category of necklaces
\newcommand{\fnec}{\mathcal{N}ec^\text{\VarFlag}\!\!} %category of flagged necklaces
\newcommand{\ffnec}{\mathcal{N}ec^\text{\VarFlag}_{f}\!\!} %category of flanked flagged necklaces
\newcommand{\ts}{S_{\otimes}} %category of ts-objects
\newcommand{\Fs}{S^{Frob}_{\otimes}} %category of Fs-objects

\newcommand{\ra}{\rightarrow} 

\newcommand{\lra}{\longrightarrow} 

\newcommand{\N}{\mathbb{N}} 

\newcommand{\uuu}{\mathcal{U}} 

\newcommand{\vvv}{\mathcal{V}} 

\newcommand{\www}{\mathcal{W}} 

\newcommand{\zzz}{\mathcal{Z}} 

\newcommand{\aaa}{\mathcal{A}} 

\newcommand{\mmm}{\mathcal{M}} 

\newcommand{\ccc}{\mathcal{C}} 

\newcommand{\ddd}{\mathcal{D}} 
\newcommand{\rrr}{\mathcal{R}} 
\newcommand{\sss}{\mathcal{S}} 

\newcommand{\op}{\mathrm{op}} 

\newcommand{\necop}{\mathcal{N}ec^{\op}}

\newcommand{\naF}{\mathsf{fF}}

\newcommand{\wK}{\mathsf{wK}}

\DeclareMathOperator{\Ext}{Ext}

\DeclareMathOperator{\Tor}{Tor}

\DeclareMathOperator{\can}{can}

\DeclareMathOperator{\nd}{nd} 

\DeclareMathOperator{\coker}{coker}

\newcommand{\CC}{\mathbf{C}} 

\DeclareMathOperator{\dayconv}{\otimes_{\Day}}

%Decorations
\newcommand{\ot}{\leftarrow}
\DeclareMathOperator{\ac}{ac}
\DeclareMathOperator{\ine}{in}
\DeclareMathOperator{\Reedy}{Reedy}
\DeclareMathOperator{\inj}{inj}
\DeclareMathOperator{\proj}{proj}
\DeclareMathOperator{\Day}{Day}

%Presentation
\newtheorem*{Ax}{Axiom}
\newtheorem{Thm}{Theorem}[section]
\newtheorem*{Thm*}{Theorem}
\newtheorem{Lem}[Thm]{Lemma}
\newtheorem{Prop}[Thm]{Proposition}
\newtheorem{Cor}[Thm]{Corollary}
\newtheorem{Q}[Thm]{Question}

\theoremstyle{definition}
\newtheorem{Def}[Thm]{Definition}
\newtheorem{Ex}[Thm]{Example}
\newtheorem{Exs}[Thm]{Examples}
\newtheorem{Con}[Thm]{Construction}
\newtheorem{Not}[Thm]{Notation}

\theoremstyle{remark}
\newtheorem{Rem}[Thm]{Remark}


\begin{document}

\begin{abstract}
In this short note, we prove that the category of necklaces is a test category. Hence presheaves on necklaces model homotopy types. The proof is analogous to that for the category of cubes with connections.
\end{abstract}

\maketitle

%\tableofcontents

\section{Introduction}

In this short note, we prove that the category $\nec$ of necklaces of \cite{dugger2011rigidification} is a test category in the sense of \cite{cisinski2006prefaisceaux}. We first recall some preliminaries on test categories and necklaces, and then prove our main result (Theorem \ref{theorem: nec is test category}).

\section{Preliminaries}

\subsection{Test categories}

We follow the terminology and conventions \cite{jardine2006categorical}. Be aware that these conventions sometimes differ from those of the original reference \cite{cisinski2006prefaisceaux}.

Let $\mathcal{A}$ be a fixed small category. We refer to presheaves $\mathcal{A}^{op}\to \Set$ as \emph{$\mathcal{A}$-sets}. Given $a\in \mathcal{A}$, let us denote the representable $\mathcal{A}$-set by $\hat{a} = \mathcal{A}(-,a)$.

Consider the functor
\begin{equation}\label{diagram: slice functor}
\iota: \mathcal{A}\rightarrow \Cat: a\mapsto \mathcal{A}/a
\end{equation}
Left Kan extension of $\iota$ along the Yoneda embedding $\mathcal{A}\hookrightarrow \Set^{\mathcal{A}^{op}}$ induces an adjunction
$$
\iota_{*} = \int_{\mathcal{A}}: \Set^{\mathcal{A}^{op}}\leftrightarrows \Cat: \iota^{*}
$$
For every $\mathcal{A}$-set $X$, $\int_{\mathcal{A}}X$ is precisely the category of elements of $X$. Given a small category $\mathcal{C}$, we have $\iota^{*}(\mathcal{C}) = \Fun(\mathcal{A}/-,\mathcal{C}): \mathcal{A}^{op}\to \Set$.

\begin{Def}
Let $N: \Cat\hookrightarrow \SSet$ denote the nerve functor.

A functor $F: \mathcal{C}\rightarrow \mathcal{D}$ between small categories is \emph{aspherical} or \emph{homotopy cofinal} if for all $d\in \mathcal{D}$, the simplicial set $N(F/d)$ is weakly contractible.  A small category $\mathcal{C}$ is \emph{aspherical} if the functor $\mathcal{C}\rightarrow *$ is aspherical, that is if $N(\mathcal{C})$ is weakly contractible.

An $\mathcal{A}$-set is called \emph{aspherical} if the forgetful functor $\int_{\mathcal{A}}X\to \mathcal{A}$ is aspherical. This is equivalent to the category $\int_{\mathcal{A}}(\hat{a}\times X)$ being aspherical for all $a\in \mathcal{A}$.
\end{Def}

\begin{Def}[Lemmas 2.3 and 2.4, \cite{jardine2006categorical}]
$\mathcal{A}$ is called a
\begin{enumerate}[1.]
\item \emph{local test category} if for every small category $\mathcal{D}$ with a terminal object, the $\mathcal{A}$-set $\iota^{*}(\mathcal{D})$ is aspherical.
\item \emph{test category} if it is a local test category and aspherical.
\end{enumerate}
\end{Def}

\subsection{Necklaces}

We follow the terminology and conventions of \cite[\S 3]{dugger2011rigidification} as well as \cite[\S 3.2]{mertens2026nerves}.

Let $\SSet_{*,*} = \SSet_{\partial\Delta^{1}/}$ denote the category of bipointed simplicial sets. Given two bipointed simplicial sets $K_{a,b}$ and $L_{c,d}$, we define their \emph{wedge sum} by the following coequalizer, which we consider to be bipointed at $a$ and $d$:
$$
\Delta^{0}\overset{b}{\underset{c}{\rightrightarrows}} K_{a,b}\amalg L_{c,d}\twoheadrightarrow K\vee L
$$


We always consider the standard simplices $\Delta^{n}$ to be bipointed at $0$ and $n$.

\begin{Def}
The category $\nec$ of \emph{necklaces} is the full subcategory of $\SSet_{*,*}$ spanned by all bipointed simplicial sets of the form
$$
T = \Delta^{n_{1}}\vee \dots\vee \Delta^{n_{k}}
$$
for $k\geq 0$ and each $n_{i} > 0$. The endpoints of each of the $\Delta^{n_{i}}$ in $T$ are called the \emph{joints} of $T$. We define the \emph{dimension} of $T$ as the number of non-joint vertices of $T$, and denote it $\dim(T)$.
\end{Def}

Under the wedge sum $\vee$, the category $\nec$ is monoidal with its unit given by $\Delta^{0}$. As a monoidal category, $\nec$ is generated by the following three classes of maps:
\begin{itemize}
\item the coface map $\delta_{j}: \Delta^{n-1}\hookrightarrow \Delta^{n}$ for $0 < j < n$,
\item the codegeneracy map $\sigma_{i}: \Delta^{n+1}\to \Delta^{n}$ for $0\leq i\leq n$, and
\item the inclusion $\nu_{k,l}: \Delta^{k}\vee \Delta^{l}\hookrightarrow \Delta^{k+l}$ which is bijective on vertices.
\end{itemize}

There is a strong monoidal functor
\begin{equation}\label{diagram: dimension functor}
\dim: \nec\to \Cat: T\mapsto [1]^{\dim(T)}
\end{equation}
where $[1] = \{0 < 1\}$ is the walking arrow category. The functor $\dim$ is defined on generating maps as follows:
\begin{itemize}
\item $\dim(\delta_{j}) = \delta^{\square}_{j,0}$
\item $\dim(\nu_{k,l}) = \delta^{\square}_{k,1}$
\item $\dim(\sigma_{i}) =
\begin{cases}
\sigma^{\square}_{1} & \text{if }i = 0\\
\gamma^{\square}_{i} & \text{if }0 < i < n\\
\sigma^{\square}_{n} & \text{if }i = n
\end{cases}$
if $n > 0$
\item $\dim(\sigma_{0}) = \id_{[0]}$ for $\sigma_{0}: \Delta^{1}\to \Delta^{0}$
\end{itemize}
where $\delta^{\square}_{i,\epsilon}$, $\sigma^{\square}_{i}$ and $\gamma^{\square}_{i}$ are maps of cubes, which are given by
\begin{align*}
\delta^{\square}_{i,\epsilon}&: [1]^{n-1}\rightarrow [1]^{n}: (\epsilon_{1},\dots,\epsilon_{n-1})\mapsto (\epsilon_{1},\dots,\epsilon_{i-1},\epsilon,\epsilon_{i},\dots,\epsilon_{n}) &&\text{for }1\leq i\leq n, \epsilon\in \{0,1\}\\
\sigma^{\square}_{i}&: [1]^{n+1}\rightarrow [1]^{n}: (\epsilon_{1},\dots,\epsilon_{n+1})\mapsto (\epsilon_{1},\dots,\epsilon_{i-1},\epsilon_{i+1},\dots,\epsilon_{n+1}) &&\text{for }1\leq i\leq n+1\\
\gamma^{\square}_{i}&: [1]^{n+1}\rightarrow [1]^{n}: (\epsilon_{1},\dots,\epsilon_{n+1})\mapsto (\epsilon_{1},\dots,\max(\epsilon_{i},\epsilon_{i+1}),\dots,\epsilon_{n}) &&\text{for }1\leq i\leq n-1
\end{align*}



\section{Proof}

We make use of the following general `yoga' as described in \cite{jardine2006categorical}.

\begin{Prop}[Remark 3.12, \cite{jardine2006categorical}]\label{proposition: amenable functor}
Let $i: \mathcal{A}\to \Cat$ be a functor and assume that
\begin{enumerate}[(a)]
\item for every $a\in \mathcal{A}$, $i(a)$ has a terminal object,
\item for every small category $\mathcal{D}$ with a terminal object, the $\mathcal{A}$-set $i^{*}(\mathcal{D})$ is aspherical, and
\item the category $\mathcal{A}$ is aspherical. 
\end{enumerate}
Then $\mathcal{A}$ is a test category.
\end{Prop}

This allows to essentially replace the functor $\iota: \mathcal{A}\to \Cat$ \eqref{diagram: slice functor} by another which may be easier to describe. In the case of necklaces, we can consider the dimension functor $\dim: \nec\to \Cat$ \eqref{diagram: dimension functor}. It induces an adjunction
$$
\dim_{*}: \Set^{\nec^{op}}\leftrightarrows \Cat: \dim^{*}
$$
where $\dim_{*}(X) = \colim_{T\in \nec}^{X_{T}}[1]^{\dim(T)}$ and $\dim^{*}(\mathcal{C}) = \Fun([1]^{\dim(-)},\mathcal{C})$.

%\begin{Lem}
%There exists an endofunctor $(-)^{*}: \nec\setminus\{\Delta^{0}\}\to \nec\setminus \{\Delta^{0}\}$ along with a commutative diagram of natural transformations:
%\[\begin{tikzcd}
%	{\Delta^{1}\vee -} && \\
%	& {(-)^{*}} & \id \\
%	\id
%	\arrow["{\nu_{1,-}}"', from=1-1, to=2-2]
%	\arrow["{\sigma_{0}\vee -}", curve={height=-6pt}, from=1-1, to=2-3]
%	\arrow["{\sigma_{0}}"', from=2-2, to=2-3]
%	\arrow["{\delta_{1}}", from=3-1, to=2-2]
%	\arrow["{=}"', curve={height=12pt}, from=3-1, to=2-3]
%\end{tikzcd}\]
%and for any necklace $U\neq\Delta^{0}$, we have $\dim(U^{*}) = \dim(U) + 1$.
%\end{Lem}
%\begin{proof}
%Given a necklace $U\neq \Delta^{0}$, we have $U = \Delta^{n}\vee U'$ for a unique necklace $U'$ and $n > 0$. Let us define $U^{*} = \Delta^{n+1}\vee U'$, so that indeed $\dim(U^{*}) = \dim(U) + 1$. Given a necklace map $f: U\to T$ define $f^{*}: U^{*}\to T^{*}$ as the unique map which is given on vertices by $f_{0}\amalg\{*\}: U^{*}_{0} = U_{0}\amalg \{*\}\to T^{*}_{0} = T_{0}\amalg \{*\}$. It is easy to see that this defines an endofunctor $(-)^{*}: \nec\to \nec$.
%
%For any necklace $U$ with $q = \Vert U\Vert > 0$, we have the necklace maps $\nu_{1,q}: \Delta^{1}\vee U\to U^{*}$, $\sigma_{0}: U^{*}\to U$ and $\delta_{1}: U\to U^{*}$ which are clearly natural in $U$ and commute as indicated above.
%\end{proof}

The next proof follows essentially the same strategy as for the category of cubes with connections, presented in \cite[Lemmas 3.10 and 3.11]{jardine2006categorical}.

\begin{Thm}\label{theorem: nec is test category}
$\nec$ is a test category.
\end{Thm}
\begin{proof}
We prove the statement by verifying conditions $(a)$-$(c)$ of Proposition \ref{proposition: amenable functor} for the functor $\dim: \nec\to \Cat$. First note that for every $T\in \nec$, the category $[1]^{\dim(T)}$ indeed has a terminal object. The category $\nec$ also has a terminal object $\Delta^{0}$, whereby it is aspherical. Hence $(a)$ and $(c)$ are satisfied.

To show $(b)$, let $\mathcal{D}$ be a small category with a terminal object $t$. We must show that for any necklace $T$, the category $\mathcal{D}_{T} = \int_{\nec}(\dim^{*}(\mathcal{D})\times \hat{T})$ is aspherical. Take an arbitrary object of $\mathcal{D}_{T}$, given by a pair $(F,f)$ with $f: U\to T$ a necklace map and $F: [1]^{\dim(U)}\rightarrow \mathcal{D}$ a functor. Further let $d = \dim(U)$ and define $f^{*} = f(\sigma_{0}\sigma_{0}\vee \id_{U}): \Delta^{2}\vee U\to U\to T$ and
$$
F^{*}: [1]^{d+1}\to \mathcal{D}: (\epsilon_{0},\dots,\epsilon_{d+1})\mapsto 
\begin{cases}
F(\epsilon_{0},\dots,\epsilon_{d}) & \text{if }\epsilon_{d+1} = 0\\
t & \text{if }\epsilon_{d+1} = 1
\end{cases}
$$
Then we have the following commutative diagrams in $\Cat$ and $\nec$ respectively:
\[\begin{tikzcd}
	{[1]^{d}} & {[1]^{d}} & {[1]^{d+1}} & {[1]^{d}} & {[1]^{d}} \\
	& & {\mathcal{D}}
	\arrow["{=}"', from=1-2, to=1-1]
	\arrow["{F}"', from=1-1, to=2-3]
	\arrow["{\delta^{\square}_{1,0}}", from=1-2, to=1-3]
	\arrow["{=}", from=1-4, to=1-5]
	\arrow["{\delta^{\square}_{1,1}}"', from=1-4, to=1-3]
	\arrow["F"{description}, from=1-2, to=2-3]
	\arrow["{F^{*}}"{description}, from=1-3, to=2-3]
	\arrow["t"{description}, from=1-4, to=2-3]
	\arrow["t", from=1-5, to=2-3]
	%\arrow["{\sigma^{\square}_{1}}"', curve={height=20pt}, dashed, from=1-3, to=1-2]
\end{tikzcd}\]
\[\begin{tikzcd}
	U & \Delta^{1}\vee U & {\Delta^{2}\vee U} & {\Delta^{1}\vee \Delta^{1}\vee U} & U \\
	& & {T}
	\arrow["\sigma_{0}\vee U"', from=1-2, to=1-1]
	\arrow["f"', from=1-1, to=2-3]
	\arrow["{\delta_{1}\vee U}", from=1-2, to=1-3]
	\arrow["{\sigma_{0}\vee \sigma_{0}\vee U}", from=1-4, to=1-5]
	\arrow["{\nu_{1,1}\vee U}"', from=1-4, to=1-3]
	\arrow["f", from=1-5, to=2-3]
	\arrow[from=1-4, to=2-3]
	\arrow[from=1-2, to=2-3]
	\arrow["f^{*}"{description}, from=1-3, to=2-3]
	%\arrow["{\sigma_{1}\vee \id_{U}}"', curve={height=15pt}, dashed, from=1-3, to=1-2]
\end{tikzcd}\]
Then we have induced endofunctors on $\mathcal{D}_{T}$:
\begin{align*}
& H: (F,f)\mapsto (F,f(\sigma_{0}\vee \id_{U})) && I: (F,f)\mapsto (F^{*},f^{*})\\
& J: (F,f)\mapsto (t,f(\sigma_{0}\vee \sigma_{0}\vee \id_{U}) && K: (F,f)\mapsto (t,f)
\end{align*}
along with natural transformations:
$$
\id_{\mathcal{D}_{T}}\leftarrow H\rightarrow I\leftarrow J\rightarrow K
$$
which induce homotopies between simplicial maps after applying the nerve. Hence, it follows that $N(K)$ is a weak homotopy equivalence. Now note that $G$ can be factored as $\mathcal{D}_{T}\xrightarrow{\alpha} \nec/T\xrightarrow{\beta} \mathcal{D}_{T}$ where $\alpha: (F,f)\mapsto f$ and $\beta: f\mapsto (t,f)$. Since clearly $\alpha\beta = \id_{\nec/T}$, it follows from the 2-out-of-6-property that $\alpha$ and $\beta$ are weak homotopy equivalence between $\mathcal{D}_{T}$ and $\nec/T$. Now $\nec/T$ is clearly aspherical and thus so is $\mathcal{D}_{T}$.
\end{proof}

\printbibliography


\end{document}
